\section{Cặp ghép}
\begingroup\small
\begin{multicols}{2}
\subsection{Cặp ghép hai phía}
Đồ thị hai phía có tập đỉnh $L,R$ và chỉ có cạnh giữa hai phía.
Cặp ghép là tập cạnh không chung đầu mút; cặp ghép cực đại theo
\emph{số cạnh} có lực lượng lớn nhất. Đừng nhầm với cặp ghép
\emph{tối đại} chỉ có nghĩa là không thêm ngay được một cạnh.
Đường tăng bắt đầu ở đỉnh chưa ghép bên $L$, luân phiên qua cạnh
chưa ghép và cạnh đã ghép, kết thúc ở đỉnh chưa ghép bên $R$.
Đảo trạng thái các cạnh trên đường này làm kích thước tăng $1$.
Một cặp ghép là cực đại khi và chỉ khi không còn đường tăng
(định lý Berge).

\subsection{Đường tăng và thuật toán Hopcroft--Karp}
Hopcroft--Karp dùng BFS tìm độ dài ngắn nhất của các đường tăng,
rồi DFS lấy nhiều đường tăng cùng độ dài trong một pha.
Độ phức tạp $O((V+E)\sqrt V)$, phù hợp khi hai phía và số cạnh lớn.
Mẫu code nhận đỉnh trái $[0,n_L)$, đỉnh phải $[0,n_R)$; gọi
\texttt{maxMatching()} một lần sau khi thêm hết cạnh.
\cppsnippet{code/matching/hopcroft_karp.hpp}

\subsection{Định lý König}
Trong đồ thị hai phía, kích thước cặp ghép cực đại bằng kích
thước phủ đỉnh cực tiểu. Sau khi tìm ghép cực đại, bắt đầu từ
các đỉnh trái chưa ghép; đi theo cạnh \emph{chưa ghép} từ trái
sang phải, rồi cạnh \emph{đã ghép} từ phải sang trái.
Gọi tập đạt được là $Z_L,Z_R$. Phủ đỉnh cực tiểu là
$(L\setminus Z_L)\cup Z_R$. Mẫu Hopcroft--Karp có hàm
\texttt{minVertexCover()} trả hai danh sách đỉnh theo công thức đó.

\subsection{Bài toán gán có trọng số}
Với $n$ người, $m\ge n$ việc, chi phí $c_{ij}$, cần gán mỗi
người một việc khác nhau để tổng chi phí nhỏ nhất. Đây là bài
toán gán \emph{đầy đủ ở phía trái}; thuật toán Hungarian chạy
$O(n^2m)$, kể cả chi phí âm. Nếu một cặp bị cấm, không tự ý
đặt chi phí vô cực rồi chấp nhận kết quả: phải kiểm tra có
gán hợp lệ hay không. Hàm dưới nhận ma trận $n\times m$ có
mọi ô hữu hạn, trả tổng chi phí và cột được gán cho từng hàng.
Trong vector đầu vào và kết quả, hàng/cột đánh số từ $0$.
Các chi phí đầu vào là \texttt{long long}, còn tổng chi phí
và các biến trung gian dùng \texttt{\_\_int128} để tránh tràn.
Bài toán tìm tổng lớn nhất có thể đổi dấu mọi chi phí rồi lấy
đối dấu kết quả, chú ý tràn với giá trị cực tiểu của kiểu số.

\vspace{-4mm}
\paragraph{Mẫu Hungarian (C++17).}
\cppsnippet{code/matching/hungarian.hpp}
\end{multicols}
\endgroup
